\chapter{系统功能详细分析与可行性}
\section{设计目标与总体定位}
声音、振动和姿态是人与物理环境交互的重要信息。单独读取一个传感器只能得到局部数值；将多源数据按共同时间组织，再连接识别与反馈功能，才能形成可理解、可操作的感知系统。本项目以板载器件为基础，把传统“读寄存器、亮LED”的实验扩展为一个可演示、可复用、可交付的多应用平台。

总体目标分为三个层次：底层实现持续且可核对的信号采集；中间层完成特征处理、事件组织和可靠传输；应用层把数据转化为实时可视化、音乐语义标签与体感控制。实时性较强的采样、时序与缓冲由 FPGA 承担，显示、数据库、音频检索和交互逻辑由电脑承担。在线识曲连接外部识别服务，构成边缘硬件与主机、网络服务协同的完整应用链路。

\fig{architecture}{多模态智能感知与交互平台总体架构}{实线为主要数据流；虚线为配置、控制与反馈。在线服务通过电脑接入，开发板自身不运行网络协议。}

本项目的智能化价值体现在对原始数据的组织与使用：声振特征在统一窗口内融合，音频录音转为歌曲名称与歌手，姿态数据转为游戏中的动作。面向 AI 感知应用，该平台具备数据采集、特征生成、任务调度和结果反馈接口；已实现的识曲采用第三方服务与本地指纹技术，后续模型训练和端侧神经网络属于可扩展方向。

\section{功能需求与完成形态}
\begin{table}[H]\centering\caption{最终系统的功能需求与实现结果}\small
\begin{tabularx}{\linewidth}{p{23mm}Y Y}\toprule
功能域 & 具体能力 & 用户可见结果\\\midrule
实时多源监测 & PCM包络、原速频谱、XYZ、两处温度、三路片内电压、开关及方向键 & 一套界面观察信号变化、有效性和运行状态。\\
声振事件记录 & 手动抓取及声音、振动、联合触发；前后触发缓存 & 保存1.5秒完整声音、XYZ及时间信息，可回放和导出。\\
桌面听歌识曲 & 三敲启动，20秒录音，在线优先与本地回退 & 显示歌曲、歌手、结果来源和阶段；支持中止与恢复就绪。\\
音乐数据管理 & 歌单导入、指纹索引、历史查询、统计与去重收藏 & 多次识别保留多条历史；同曲重复收藏保持一份。\\
板载反馈 & 数码管、16路LED、RGB与SW15模式开关 & 不持续盯着电脑也能了解就绪、录音、传输及结果状态。\\
人机交互应用 & 本地抢歌名；八关《重力迷航》 & 比较抢答时刻；通过倾斜与中键控制游戏，也可键盘练习。\\
工程交付 & 独立EXE、设备发现、软件内初始化、诊断与校验 & 无需开发环境即可运行，源码和版本可追溯。\\\bottomrule
\end{tabularx}\end{table}

三个应用共享感知基础设施，但按任务需求使用不同的数据产品。实时界面使用轻量预览；事件模块使用带校验的完整快照；识曲模块使用独立长录音；游戏使用持续更新的 XYZ 与按键。将这些路径分别设计，避免把“画面刷新”误当作“完整数据保存”，也使新增应用能够复用同一硬件服务。

\section{典型使用场景与交互闭环}
在桌面监测场景中，麦克风感知空气传播的声音，加速度传感器感知通过桌面传到板上的振动。用户可比较两类信号并保存触发前后的上下文。事件文件同时包含声音、XYZ、触发原因与参数版本，能够回答“发生了什么”和“系统为何保存”。

在音乐场景中，手机外放作为声源，三次连续敲击成为自然交互入口。系统确认就绪后采集音乐，使用灯光和数码管反馈进度，再由主机给出歌曲标签。结果、原始录音和来源信息一起保存，BTNC 中键或界面按钮完成收藏。相比单纯调用识曲接口，工程还覆盖了物理唤醒、完整采集、跨设备传输、失败回退与结果管理。

在体感场景中，开发板被作为可握持的控制器。游戏先估计自然姿态并校准方向，再把相对倾斜转为运动输入。按住中键可刹车，页面切换和失焦触发暂停。八关地图、机关、星级与存档形成持续体验，使传感采集的成果能够直接转化为人机交互。

\section{硬件、资源与成本可行性}
Nexys A7-100T 已集成 Artix-7 FPGA、数字麦克风、ADXL362 加速度计、ADT7420 温度传感器、USB通信与显示输入资源\cite{rm}。核心功能无需新增外部传感器；电脑承担界面与网络服务，降低了板端算法和存储压力。主要投入是已有开发板、USB数据线、Windows电脑及播放音乐的手机。

事件声音的数据量为
\begin{equation}
 S_\mathrm{event}=16000\times1.5\times2=48000\ \mathrm{B}.
\end{equation}
按标称200组/秒、每组10字节估算，XYZ每事件约3 kB；双组事件数据约102 kB。音乐录音采用独立缓冲：
\begin{equation}
 S_\mathrm{music}=8000\times20\times2=320000\ \mathrm{B}.
\end{equation}
将20秒录音降到8 kHz后，仍保持160000点容量，与10秒16 kHz录音共用同一存储上限。最终 BRAM 使用率为88.52\%，资源较紧但已完成实现；第4章进一步分析降采样和存储组织。

UART采用460800 bit/s、8N1。每字节占10个线路位，理想有效吞吐约46080 B/s，因此320000字节录音仅数据传输就至少需要6.94秒，实际还包括请求、响应、校验和调度开销。系统将录音、传输和识别作为不同阶段显示，允许用户理解等待来源，不将服务返回时间当作全过程时间。

\section{软件与技术路线可行性}
主机以 Python、Tk/ttkbootstrap、NumPy、Matplotlib 和 SQLite 组成。它们分别承担界面、数值处理、可视化与本地持久化；单文件打包将运行环境、配套固件和必要组件一并交付。FPGA 使用已有 Vivado 工具链与官方器件约束，设计由顶层模块、独立接口控制器和可验证的处理核心组成。

在线识曲通过 ShazamIO 调用服务，本地备用采用 audfprint 的频谱地标指纹检索\cite{shazamio,audfprint}。本地方案不要求训练样本即可建立有限歌单；在线方案借助外部音乐目录覆盖更多曲目。两种路径共享同一完整录音，结果明确显示来源。网络异常时保留音频并解释原因，避免系统功能完全依赖单次网络成功。

\begin{table}[H]\centering\caption{需求约束与对应设计选择}\small
\begin{tabularx}{\linewidth}{p{28mm}Y}\toprule 约束 & 设计选择\\\midrule
多源采样频率不同 & 单一100 MHz逻辑时钟，使用使能脉冲与实际采样时间戳组织数据。\\
BRAM容量有限 & 事件与识曲分开缓存，20秒音乐使用抗混叠降采样，应用模式协同调度。\\
串口带宽有限 & 实时极值预览与完整快照分离，按编号分块读取，显示实际传输阶段。\\
界面需要保持响应 & 串口服务、GUI与识曲进程分工；主线程通过队列更新界面。\\
现场环境可能变化 & 阈值可配置、游戏先校准、异常保留诊断；标定边界集中说明。\\\bottomrule
\end{tabularx}\end{table}

通过上述分工，项目将“实现某一个功能”扩展为“在资源与时序约束下组织多个应用”。功能分析不仅列出界面按钮，也明确每项功能的数据来源、资源需求和可观察验收结果。
